\باب{تحدیدی مسائل   کے ثبوت}\شناخت{ضمیمہ_ب}


اس ضمیمے میں حصہ \حوالہ{حصہ_حد_قواعد} میں پیش کیے گئے  مسئلہ \حوالہ{مسئلہ_حد_قواعد-الف} کے جزو  \اردوعدد{2} تا \اردوعدد{5}  اور مسئلہ \حوالہ{مسئلہ_حد_بیچ}کے   ثبوت  دیے جاتے ہیں۔

پہلے  حصہ \حوالہ{حصہ_حد_قواعد}  کا درج ذیل    مسئلہ \حوالہ{مسئلہ_حد_قواعد-الف}  دیکھتے ہیں۔
\ابتدا{مسئلہ}\شناخت{مسئلہ_ضمیمہ_دو_ایک}\موٹا{حد کے خواص}\\
اگر \عددی{L}، \عددی{M}، \عددی{c}، اور \عددی{k} حقیقی اعداد ہوں اور
\[\lim_{x\to c} f(x)=L,\quad \lim_{x\to c}g(x)=M\]
ہوں تب درج ذیل ہوں گے۔
\begin{center}
\begin{tabular}{rlrl}
\اردوعدد{1}&مجموعے کا قاعدہ&\(\lim_{x\to c} (f(x)+g(x))=L+M\) & \\
\اردوعدد{2}&فرق  کا قاعدہ&\(\lim_{x\to c} (f(x)-g(x))=L-M\)& \\
\اردوعدد{3}&حاصل ضرب کا قاعدہ&\(\lim_{x\to c} (f(x)\cdot g(x))=L\cdot M\) &\\
\اردوعدد{4}&مستقل سے ضرب کا قاعدہ&\(\lim_{x\to c} (kf(x))=kL\) & ($k$ کوئی بھی عدد ہو سکتا ہے)\\
\اردوعدد{5}&حاصل تقسیم  کا قاعدہ&\(\lim_{x\to c}\frac{f(x)}{g(x)} =\frac{L}{M}\) & جہاں \عددی{M\ne 0} ہے \\
\اردوعدد{6}&طاقت   کا قاعدہ&
\multicolumn{2}{l}{\multirow{1}{*}{
اگر \عددی{r} اور \عددی{s} عدد صحیح ہوں  جن کے مشترک جزو  ضربی موجود نہ ہوں اور \عددی{s\ne 0} ہو تب}}\\
&&\parbox[c][1.75cm][c]{4cm}
{\centering
\(\lim_{x\to c} (f(x))^{r/s}=L^{r/s}\) }
&\\
&&
\multicolumn{2}{l}{\multirow{1}{*}{ہو گا بشرطیکہ   \عددی{L^{r/s}} حقیقی عدد ہو۔(جفت  \عددی{s} کی  صورت میں ہم \عددی{L>0} فرض کرتے ہیں۔)}}
\end{tabular}
\end{center}
\انتہا{مسئلہ}

ہم نے مجموعے کے قاعدے  کا ثبوت حصہ \حوالہ{حصہ_حد_مطلوبہ_قیمتیں_اور_حد}  میں پیش کیا جبکہ  طاقتی قاعدے کا ثبوت  اعلٰی نسابی کتابوں میں  ملے گا۔مجموعے کے قاعدہ میں  \عددی{g(x)} کی جگہ \عددی{-g(x)} اور \عددی{M} کی جگہ \عددی{-M} ڈالنے سے  فرق کا قاعدہ حاصل ہو گا۔ حاصل ضرب قاعدہ میں \عددی{g(x)=k} ڈالنے سے مستقل ضرب کا قاعدہ حاصل ہو گا۔ یوں صرف حاصل ضرب اور حاصل  تقسیم کے قاعدے  رہ گئے ہیں۔


\ابتدا{ثبوت}\موٹا{تحدیدی حاصل ضرب  قاعدے کا ثبوت}\quad 
ہم دکھاتے ہیں کہ  کسی بھی  \عددی{\epsilon>0} کے لئے ایسا \عددی{\delta>0} موجود ہو گا  کہ    \عددی{f}  اور \عددی{g}  کے دائرہ کار کے تقاطع
  \عددی{D}  میں تمام \عددی{x}  پر درج ذیل ہو گا۔
\[0<|x-c|<\delta \qquad \Rightarrow \qquad   |f(x)g(x)-LM|<\epsilon\]
اب فرض کریں   \عددی{\epsilon} مثبت عدد ہے اور \عددی{f(x)} اور \عددی{g(x)} کو درج ذیل  لکھیں۔
\[f(x)=L+(f(x)-L),\qquad g(x)=M+(g(x)-M)\]
ان کو آپس میں ضرب دے کر \عددی{LM} منفی کرنے سے درج ذیل ملتا ہے۔
\begin{align}
f(x)\cdot g(x)-LM&=(L+(f(x)-L))(M+(g(x)-M))-LM\notag\\
&=LM+L(g(x)-M)+M(f(x)-L)\notag\\
&\quad\quad\quad\quad+(f(x)-L)(g(x)-M)-LM\notag\\
&=L(g(x)-M)+M(f(x)-L)+(f(x)-L)(g(x)-M)\label{مساوات_ضمیمہ_ب_ضرب_ثبوت}
\end{align}
چونکہ \عددی{x\to c}  کرنے سے \عددی{f} اور \عددی{g} کے حد \عددی{L} اور \عددی{M}  پائے جاتے ہیں، لہٰذا \عددی{D} میں تمام \عددی{x} کے لئے ایسے مثبت  اعداد \عددی{\delta_1}، \عددی{\delta_2}، \عددی{\delta_3}، اور \عددی{\delta_4} موجود ہوں گے کہ درج ذیل ہوں گے۔
\begin{equation}
\begin{aligned}\label{مساوات_ضمیمہ_الف_چار}
&0<|x-c|<\delta_1 \qquad \Rightarrow  \qquad |f(x)-L|<\sqrt{\epsilon/3}\\
&0<|x-c|<\delta_2 \qquad \Rightarrow  \qquad |f(x)-L|<\sqrt{\epsilon/3}\\
&0<|x-c|<\delta_3 \qquad \Rightarrow  \qquad |f(x)-L|<\epsilon/(3(1+|M|))\\
&0<|x-c|<\delta_4 \qquad \Rightarrow  \qquad |f(x)-L|<\epsilon/(3(1+|L|))
\end{aligned}
\end{equation}
اگر ہم \عددی{\delta_1} تا \عددی{\delta_4} میں سب سے چھوٹے کو \عددی{\delta} لیں ، عدم 
مساوات \حوالہ{مساوات_ضمیمہ_الف_چار}  کے دائیں ہاتھ کی  عدم مساوات \عددی{0<|x-c|<\delta}کے لئے بھی مطمئن ہوں گی۔ لہٰذا  \عددی{D} میں تمام \عددی{x} کے لئے ،    \عددی{0<|x-c|<\delta}  سے مراد  درج ذیل لیا جا سکتا ہے۔
\begin{align*}
|f(x)\cdot  & g(x)-LM|\quad \quad\quad \quad \quad \text{\RL{مساوات \حوالہ{مساوات_ضمیمہ_ب_ضرب_ثبوت} پر مثلثی عدم مساوات کا اطلاق}}\\
&\le |L||g(x)-M|+|M||f(x)-L|+|f(x)-L||g(x)-M|\\
&\le (1+|L|)|g(x)-M|+(1+|M|)|f(x)-L|+|f(x)-L||g(x)-M|\\
&<\frac{\epsilon}{3}+\frac{\epsilon}{3}+\sqrt{\frac{\epsilon}{3}}\sqrt{\frac{\epsilon}{3}}=\epsilon \quad \quad
 \text{\RL{مساوات \حوالہ{مساوات_ضمیمہ_الف_چار}سے قیمتیں لی گئیں}}
\end{align*}
یوں تحدیدی حاصل ضرب قاعدے  کا ثبوت مکمل ہوتا ہے۔
\انتہا{ثبوت}

\ابتدا{ثبوت}\موٹا{تحدیدی حاصل تقسیم قاعدے کا ثبوت}\quad 
ہم دکھاتے ہیں کہ \عددی{\lim_{x\to c} (1/g(x))=1\!/\!M} ہو گا۔ اس کے بعد ہم  تحدیدی حاصل ضرب قاعدہ کے تحت   درج ذیل  اخذ کر سکتے ہیں۔
\[\lim_{x\to c}\frac{f(x)}{g(x)}=\lim_{x\to c}\left(f(x)\cdot \frac{1}{g(x)}\right)=\lim_{x\to c}f(x)\cdot \lim_{x\to c}\frac{1}{g(x)}=L\cdot \frac{1}{M}=\frac{l}{M}\]

فرض کریں \عددی{\epsilon>0} ہے۔ یہ دکھانے کی خاطر کہ  \عددی{\lim_{x\to c} (1/g(x))=1\!/\!M} ہے ہمیں دکھانا ہو گا کہ تمام \عددی{x} کے لئے ایسا \عددی{\delta>0} موجود ہے کہ درج ذیل مطمئن ہو۔
\[0<|x-c|<\delta \quad \Rightarrow \quad \left | \frac{1}{g(x)}-\frac{1}{M}\right |<\epsilon\]

چونکہ \عددی{|M|>0} ہے لہٰذا تمام \عددی{x} کے لئے ایسا مثبت عدد \عددی{\delta_1} موجود ہے  جو درج ذیل کو مطمئن کرتا ہے۔
\begin{align}\label{مساوات_ضمیمہ_ب_تعلق}
0<|x-c|<\delta_1\quad \Rightarrow \quad |g(x)-M|<\frac{M}{2}
\end{align}
کسی بھی عدد \عددی{A} اور \عددی{B} کے لئے  ہم جانتے ہیں کہ \عددی{|A|-|B|\le |A-B|} اور \عددی{|B|-|A|\le |A-B|} ہو گا جس 
سے \عددی{||A|-|B||\le |A-B|} اخذ کیا جا سکتا ہے۔ اس میں  \عددی{A=g(x)} اور \عددی{B=M} لینے سے درج ذیل  حاصل ہو گا
\[||g(x)|-|M||\le |g(x)-M|\]
 جس کو عدم مساوات \حوالہ{مساوات_ضمیمہ_ب_تعلق} کے ساتھ ملا کر درج ذیل لکھا جا سکتا ہے۔
\begin{align}
||g(x)|-|M||<\frac{|M|}{2}\notag\\
-\frac{|M|}{2}<|g(x)|-|M|<\frac{|M|}{2}\notag\\
\frac{|M|}{2}<|g(x)|<\frac{3|M|}{2}\notag\\
|M|<2|g(x)|<3|M|\notag\\
\frac{1}{|g(x)|}<\frac{2}{|M|}<\frac{3}{|g(x)|}\label{مساوات_ضمیمہ_ب_چار}
\end{align}
یوں \عددی{0<|x-c|<\delta_1} سے درج ذیل مراد لیا جا سکتا ہے۔
\begin{align}\label{مساوات_ضمیمہ_ب_پانچ}
\left | \frac{1}{g(x)}-\frac{1}{M}\right |&=\left | \frac{M-g(x)}{Mg(x)}\right |\le \frac{1}{|M|}\cdot \frac{1}{|g(x)|}\cdot |M-g(x)|  \notag\\
&<\frac{1}{|M|}\cdot \frac{2}{|M|}\cdot |M-g(x)|\quad \text{\RL{عدم مساوات \حوالہ{مساوات_ضمیمہ_ب_چار}}}
\end{align}
چونکہ \عددی{(1/2)|M|^2\epsilon>0} ہے لہٰذا تمام \عددی{x} کے لئے ایسا عدد \عددی{\delta_2>0} موجود ہو گا جو درج ذیل کو مطمئن کرتا ہو۔
\begin{align}\label{مساوات_ضمیمہ_ب_چھ}
0<|x-c|<\delta_2\quad \Rightarrow \quad |M-g(x)|<\frac{\epsilon}{2}|M|^2
\end{align}
اگر ہم \عددی{\delta} کو \عددی{\delta_1} اور \عددی{\delta_2} میں چھوٹے کے برابر لیں، تب  تمام \عددی{x} کے لئے  \عددی{0<|x-c|<\delta} کی صورت میں عدم مساوات \حوالہ{مساوات_ضمیمہ_ب_پانچ} اور \حوالہ{مساوات_ضمیمہ_ب_چھ} مطمئن ہوں گے۔ ان نتائج کو ملا کر درج ذیل حاصل ہو گا۔
\[0<|x-c|<\delta \quad \Rightarrow \quad \left|\frac{1}{g(x)}-\frac{1}{M}\right |<\epsilon \]
یوں تحدیدی حاصل تقسیم کے قاعدے کا ثبوت مکمل  ہوتا ہے۔
\انتہا{ثبوت}

اب حصہ \حوالہ{حصہ_حد_قواعد}  کا درج ذیل    مسئلہ \حوالہ{مسئلہ_حد_بیچ}  دیکھتے ہیں۔

\ابتدا{مسئلہء}\موٹا{مسئلہ بیچ}\quad 
فرض کریں کھلے  وقفہ \عددی{I}  میں، جس میں \عددی{c} پایا جاتا ہو ،  (غالباً) \عددی{x=c} کے سوا تمام \عددی{x} کے
 لئے درج ذیل ہے۔
\[{g(x)\le f(x)\le h(x)}\]
 مزید فرض کریں \عددی{\lim_{x\to c}g(x)=\lim_{x\to c}h(x)=L} ہے۔ تب \عددی{\lim_{x\to c}f(x)=L} ہو گا۔
\انتہا{مسئلہء}

\ابتدا{ثبوت}\موٹا{دائیں ہاتھ حدوں  کا ثبوت}\quad
فرض کریں  \عددی{\lim_{x\to c^+}g(x)=\lim_{x\to c^+}h(x)=L} ہے۔ تب  کسی بھی \عددی{{\epsilon>0}} کے لئے ایسا \عددی{\delta>0} موجود ہو گا کہ تمام \عددی{x} کے لئے وقفہ \عددی{c<x<c+\delta} وقفہ \عددی{I}   کے اندر واقع ہو اور  عدم مساوات سے مراد درج ذیل ہو۔
\[L-\epsilon<g(x)<L+\epsilon\quad \text{اور} \quad L-\epsilon<h(x)<L+\epsilon\]
ان عدم مساوات  کو عدم مساوات  \عددی{g(x)\le f(x)\le h(x)}    کے ساتھ ملا کر درج ذیل ہو گا۔
\begin{align*}
L-\epsilon&<g(x)\le f(x)\le h(x)<L+\epsilon,\\
L-\epsilon&<f(x)<L+\epsilon,\\
-\epsilon&<f(x)-L<\epsilon
\end{align*}
یوں تمام \عددی{x} کے لئے عدم مساوات \عددی{c<x<c+\delta} سے مراد \عددی{|f(x)-L|<\epsilon} ہو گا۔
\انتہا{ثبوت}

\ابتدا{ثبوت}\موٹا{بائیں ہاتھ حدوں  کا ثبوت}\quad
فرض کریں \عددی{\lim_{x\to c^-}g(x)=\lim_{x\to c^-}h(x)=L} ہے۔ تب کسی بھی \عددی{{\epsilon>0}} کے لئے  ایسا \عددی{\delta>0} موجود ہو گا کہ تمام \عددی{x}  کے لئے وقفہ \عددی{c-\delta<x<c} وقفہ \عددی{I} کے اندر واقع ہو اور عدم مساوات سے درج ذیل مراد ہو۔
\[L-\epsilon<g(x)<L+\epsilon\quad\text{اور}\quad L-\epsilon<h(x)<L+\epsilon\]
پہلے کی طرح ہم یہ نتیجہ اخذ کرتے ہیں کہ  تمام \عددی{x} کے لئے عدم مساوات \عددی{c-\delta<x<c} سے مراد \عددی{|f(x)-L|<\epsilon} ہو گا۔
\انتہا{ثبوت}

\ابتدا{ثبوت}\موٹا{دو طرفہ حدوں کا ثبوت}\quad 
اگر \عددی{\lim_{x\to c}g(x)=\lim_{x\to c}h(x)=L} ہو، تب  \عددی{x\to c^+} اور \عددی{x\to c^-} دونوں  کے لئے \عددی{g(x)} اور \عددی{h(x)} حد \عددی{L} کو پہنچتے ہیں۔ یوں \عددی{\lim_{x\to c^+}f(x)=L} اور  \عددی{\lim_{x\to c^-}f(x)=L}  ہوں گے۔ لہٰذا \عددی{\lim_{x\to c}f(x)} موجود   اور \عددی{L} کے برابر ہو گا۔
\انتہا{ثبوت}

%===============================================
\حصہء{سوالات}

\ابتدا{سوال}\شناخت{سوال_ضمیمہ_دو_ایک}
فرض کریں  \عددی{x\to c}  کرنے سے تفاعل \عددی{f_1(x)}، \عددی{f_2(x)}، اور \عددی{f_3(x)} کے حد بالترتیب \عددی{L_1}، \عددی{L_2}، اور \عددی{L_3} ہیں۔دکھائیں کہ ان تفاعلات کے مجموعے کا حد \عددی{L_1+L_2+L_3} ہو گا۔ ریاضیاتی ترغیب (ضمیمہ الف) استعمال کر کے اس نتیجے کو  محدود تعداد کے تفاعلات کے مجموعے تک عمومیت دو۔
\انتہا{سوال}
\ابتدا{سوال}\شناخت{سوال_ضمیمہ_دو_دو}
مسئلہ \حوالہ{حصہ_حد_قواعد}   میں حاصل ضرب کی  حد   اور ریاضیاتی  ترغیب استعمال  کر کے دکھائیں کہ اگر تفاعل \عددی{f_1(x)}،   \عددی{f_2(x)}، \نقطے،  \عددی{f_n(x)} کے حد بالترتیب \عددی{L_1}،  \عددی{L_2}،\نقطے، \عددی{L_n} ہوں تب درج ذیل ہو گا۔
\[\lim_{x\to c} f_1(x)\lim_{x\to c} f_2(x)\cdots \lim_{x\to c} f_n(x)=L_1 L_2\cdots L_n\]
\انتہا{سوال}
\ابتدا{سوال}\شناخت{سوال_ضمیمہ_دو_تین}
یہ جانتے ہوئے کہ \عددی{\lim_{x\to c}x=c} ہے اور سوال \حوالہ{سوال_ضمیمہ_دو_دو} کا نتیجہ استعمال  کرکے دکھائیں کہ کسی بھی عدد صحیح \عددی{n>1} کے لئے \عددی{\lim_{x\to c}x^n=c^n} ہو گا۔
\انتہا{سوال}
%Q4
\ابتدا{سوال}\شناخت{سوال_ضمیمہ_دو_چار}\موٹا{کثیر رکنی کی حد}\quad 
یہ جانتے ہوئے کہ کسی بھی عدد \عددی{k} کے لئے \عددی{\lim_{c\to c}(k)=k} ہو گا اور سوال \حوالہ{سوال_ضمیمہ_دو_ایک} اور سوال \حوالہ{سوال_ضمیمہ_دو_تین} کے نتائج استعمال کر کے دکھائیں کہ  کسی بھی کثیر رکنی تفاعل
\[f(x)=a_{n}x^{n}+a_{n-1}x^{n-1}+\cdots+a_{1}x^{1}+a_{0}\]
کے لئے \عددی{\lim_{x\to c}f(x)=f(c)} ہو گا۔
\انتہا{سوال}
\ابتدا{سوال}\موٹا{ناطق تفاعلات کی حدیں}\quad 
مسئلہ \حوالہ{مسئلہ_ضمیمہ_دو_ایک}  اور سوال \حوالہ{سوال_ضمیمہ_دو_چار}  کا نتیجہ استعمال کرتے ہوئے دکھائیں کہ  اگر  \عددی{f(x)} اور \عددی{g(x)} کثیر رکنی تفاعلات  ہوں اور \عددی{g(c)\ne 0} ہو تب درج ذیل ہو گا۔
\[\lim_{x\to c}\frac{f(x)}{g(x)}=\frac{f(c)}{g(c)}\]
\انتہا{سوال}
\ابتدا{سوال}\موٹا{استمراری  تفاعلات کے مرکب}\quad 
  دو استمراری تفاعلات کا  مرکب  استمراری ہونے  کا خاکہ شکل  \حوالہ{شکل_ضمیمہ_ایک_ایک}  میں پیش کیا گیا ہے۔ خاکے سے ثبوت کی   تشکیل کریں۔آپ نے ثابت کرنا ہے کہ  اگر \عددی{x=c} پر \عددی{f} استمراری ہو اور \عددی{f(c)} پر \عددی{g} استمراری ہو تب \عددی{c} پر \عددی{g\circ  f} استمراری ہو گا۔

\begin{figure}
\centering
\begin{tikzpicture}
\draw[thick](0,0)--++(2,0)node[pos=0.5,circ]{}node[pos=0.5,below]{$c$};
\draw[thick](3,0)--++(2,0)node[pos=0.5,circ]{}node[pos=0.5,below]{$f(c)$};
\draw[thick](6,0)--++(2,0)node[pos=0.5,circ]{}node[pos=0.5,below]{$g(f(c))$};
\draw[thick,-stealth](1.75,0.4) to [out=30,in=150]node[pos=0.65,above]{$f$}(3.25,0.4);
\draw[thick,-stealth](4.75,0.4) to [out=30,in=150]node[pos=0.4,above]{$g$}(6.25,0.4);
\draw[thick,-stealth](1.5,0.4) to [out=45,in=135]node[pos=0.5,above]{$g\circ f$}(6.5,0.4);
\draw[](0.25,0)node[]{$($} (1.75,0)node[]{$)$};
\draw[](3.25,0)node[]{$($} (4.75,0)node[]{$)$};
\draw[](6.25,0)node[]{$($} (7.75,0)node[]{$)$};
\draw [decorate,decoration={brace,amplitude=3pt,raise=1.5ex}]
  (0.25,0) -- (0.98,0) node[midway,yshift=1.75em]{$\delta_f$};
\draw [decorate,decoration={brace,amplitude=3pt,raise=1.5ex}]
  (1.01,0) -- (1.75,0) node[midway,yshift=1.75em]{$\delta_f$};

\draw [decorate,decoration={brace,amplitude=3pt,raise=1.5ex}]
  (3.25,0) -- (3.98,0) node[midway,yshift=1.75em]{$\delta_g$};
\draw [decorate,decoration={brace,amplitude=3pt,raise=1.5ex}]
  (4.01,0) -- (4.75,0) node[midway,yshift=1.75em]{$\delta_g$};

\draw [decorate,decoration={brace,amplitude=3pt,raise=1.5ex}]
  (6.25,0) -- (6.98,0) node[midway,yshift=1.75em]{$\epsilon$};
\draw [decorate,decoration={brace,amplitude=3pt,raise=1.5ex}]
  (7.01,0) -- (7.75,0) node[midway,yshift=1.75em]{$\epsilon$};
\end{tikzpicture}
\caption{دو استمراری تفاعلات کا  مرکب بھی استمراری ہو گا۔}
\label{شکل_ضمیمہ_ایک_ایک}
\end{figure}
فرض کریں کہ   نقطہ \عددی{c}   تفاعل \عددی{f} کے دائرہ کار  کے اندر  واقع ہے اور  نقطہ \عددی{f(c)} تفاعل \عددی{g} کے دائرہ کار کے اندر  واقع  ہے۔ ایسا کرنے سے مطلوبہ  حدیں دو طرفہ صورت اختیار  کرتی ہیں۔ (یک طرفہ حد کی صورت میں بھی  دلائل   اسی طرح  ہیں۔)
\انتہا{سوال}

